\section{Hạn chế của đề tài}

Bên cạnh những kết quả đạt được, đề tài có một số hạn chế cần ghi nhận:

\begin{itemize}
    \item \textbf{Dữ liệu lỗi giả lập:} Do các sự cố hỏng hóc tự nhiên hiếm gặp, dữ liệu bất thường chủ yếu sử dụng phương pháp giả lập lỗi. Điều này chưa phản ánh sự thoái hóa cơ khí tuyến tính như mòn vòng bi qua nhiều năm.
    
    \item \textbf{Ngưỡng cắt cứng (Hard-thresholding):} Sử dụng ngưỡng phương sai cố định để phân pha ba trạng thái có thể gây dao động trạng thái (hiệu ứng ping-pong) trong giai đoạn chuyển đổi giữa các chu kỳ giặt.
    
    \item \textbf{Sự phụ thuộc vào firmware tĩnh:} Hệ thống chưa có khả năng tự học và thích ứng (Học trên thiết bị) với sự thay đổi thiết bị theo thời gian. Cần cập nhật firmware từ xa khi các thông số cơ khí thay đổi đáng kể.

    \item \textbf{Chưa đánh giá tiêu thụ năng lượng dài hạn:} Đồ án đã đo đạc độ trễ suy luận, dung lượng Flash và khả năng vận hành thời gian thực, nhưng chưa thực hiện phép đo dòng tiêu thụ hoặc công suất trung bình trong các phiên chạy nhiều ngày. Vì vậy, khả năng triển khai dưới dạng nút cảm biến chạy pin vẫn cần được kiểm chứng thêm.
\end{itemize}
